\section{Spacetime Interval and Metric}
\label{sec:spacetime_interval_and_metric}

Different inertial observers assign different coordinate differences to the same pair of events. A relativistic geometry therefore requires a quantity that combines temporal and spatial separations and has the same value in every inertial frame.

That quantity is the spacetime interval. It is defined using the Minkowski metric, which also establishes the sign convention used throughout the remainder of these notes.

\subsection{The metric convention}

We adopt the mostly-minus signature

\[
\boxed{
\eta_{\mu\nu}
=
\operatorname{diag}(1,-1,-1,-1)
}.
\]

In matrix form,

\[
\eta_{\mu\nu}
=
\begin{pmatrix}
1&0&0&0\\
0&-1&0&0\\
0&0&-1&0\\
0&0&0&-1
\end{pmatrix}.
\]

The symbol \(\eta_{\mu\nu}\) denotes the metric of flat Minkowski spacetime. Greek indices run over

\[
\mu,\nu=0,1,2,3.
\]

The temporal component carries a positive sign, while each spatial component carries a negative sign.

This convention must be used consistently. A text using the opposite signature

\[
(-,+,+,+)
\]

will display sign changes in norms, raised derivatives, the d'Alembert operator, and some Lagrangian densities, even though the physical predictions are unchanged.

\subsection{The spacetime interval}

For an infinitesimal displacement

\[
\mathrm{d}x^{\mu}
=
(c\,\mathrm{d}t,\mathrm{d}x,\mathrm{d}y,\mathrm{d}z),
\]

the spacetime interval is

\[
\boxed{
\mathrm{d}s^{2}
=
\eta_{\mu\nu}
\mathrm{d}x^{\mu}
\mathrm{d}x^{\nu}
}.
\]

Expanding the components gives

\[
\boxed{
\mathrm{d}s^{2}
=
c^{2}\mathrm{d}t^{2}
-
\mathrm{d}x^{2}
-
\mathrm{d}y^{2}
-
\mathrm{d}z^{2}
}.
\]

Using

\[
\mathrm{d}\mathbf{x}^{2}
=
\mathrm{d}x^{2}
+
\mathrm{d}y^{2}
+
\mathrm{d}z^{2},
\]

we may write

\[
\mathrm{d}s^{2}
=
c^{2}\mathrm{d}t^{2}
-
\mathrm{d}\mathbf{x}^{2}.
\]

For two finite events, define

\[
\Delta x^{\mu}
=
x_{2}^{\mu}-x_{1}^{\mu}.
\]

Their squared separation is

\[
\boxed{
\Delta s^{2}
=
\eta_{\mu\nu}
\Delta x^{\mu}
\Delta x^{\nu}
=
c^{2}(\Delta t)^{2}
-
|\Delta\mathbf{x}|^{2}
}.
\]

\subsection{Why the interval is not an ordinary Euclidean distance}

In Euclidean space, the squared distance between nearby points is

\[
\mathrm{d}\ell^{2}
=
\mathrm{d}x^{2}
+
\mathrm{d}y^{2}
+
\mathrm{d}z^{2}.
\]

Every nonzero Euclidean squared distance is positive. In Minkowski spacetime, by contrast,

\[
\mathrm{d}s^{2}
=
c^{2}\mathrm{d}t^{2}-\mathrm{d}\mathbf{x}^{2}
\]

may be positive, zero, or negative.

This indefinite sign is not a technical inconvenience. It encodes causal structure. Temporal and spatial directions are geometrically different, and the sign of \(\Delta s^{2}\) determines which events can be causally connected.

\subsection{Timelike, null, and spacelike separations}

With the chosen signature, the three cases are:

\begin{enumerate}
    \item Timelike separation:
    \[
    \boxed{\Delta s^{2}>0}.
    \]
    Then
    \[
    c^{2}(\Delta t)^{2}>|\Delta\mathbf{x}|^{2}.
    \]
    A massive particle moving slower than light can connect the events.

    \item Null or lightlike separation:
    \[
    \boxed{\Delta s^{2}=0}.
    \]
    Then
    \[
    |\Delta\mathbf{x}|=c|\Delta t|.
    \]
    A light signal can connect the events.

    \item Spacelike separation:
    \[
    \boxed{\Delta s^{2}<0}.
    \]
    Then
    \[
    |\Delta\mathbf{x}|>c|\Delta t|.
    \]
    Connecting the events would require a signal faster than light.
\end{enumerate}

The sign assignments reverse when the opposite metric signature is used. The causal classification itself does not change.

\subsection{Lorentz invariance of the interval}

A Lorentz transformation is a linear coordinate transformation

\[
x'^{\mu}
=
\Lambda^{\mu}{}_{\nu}x^{\nu}
\]

that preserves the Minkowski metric:

\[
\boxed{
\Lambda^{\mathsf{T}}\eta\Lambda
=
\eta
}.
\]

Therefore

\[
\begin{aligned}
\Delta s'^{2}
&=
\eta_{\mu\nu}
\Delta x'^{\mu}
\Delta x'^{\nu}\\
&=
\eta_{\mu\nu}
\Lambda^{\mu}{}_{\alpha}
\Lambda^{\nu}{}_{\beta}
\Delta x^{\alpha}
\Delta x^{\beta}\\
&=
\eta_{\alpha\beta}
\Delta x^{\alpha}
\Delta x^{\beta}\\
&=
\Delta s^{2}.
\end{aligned}
\]

Thus

\[
\boxed{
\Delta s'^{2}=\Delta s^{2}
}.
\]

The explicit form of \(\Lambda^{\mu}{}_{\nu}\) and a direct component verification will be given later. Here the defining geometric property is already visible: Lorentz transformations preserve the interval.

\subsection{Lowering the coordinate index}

The metric converts a contravariant coordinate into a covariant coordinate:

\[
x_{\mu}
=
\eta_{\mu\nu}x^{\nu}.
\]

For

\[
x^{\mu}=(ct,x,y,z),
\]

we obtain

\[
\boxed{
x_{\mu}=(ct,-x,-y,-z)
}.
\]

The interval may then be written compactly as

\[
\Delta s^{2}
=
\Delta x_{\mu}\Delta x^{\mu}.
\]

This is a contraction: the repeated index appears once upstairs and once downstairs and is summed over all four values.

The procedure of raising and lowering general vector indices will be developed fully in the next section.

\subsection{Proper time}

For a timelike worldline, the proper time \(\tau\) is defined by

\[
\boxed{
c^{2}\mathrm{d}\tau^{2}
=
\mathrm{d}s^{2}
}.
\]

Using

\[
\mathrm{d}\mathbf{x}
=
\mathbf{v}\,\mathrm{d}t,
\]

we find

\[
\begin{aligned}
c^{2}\mathrm{d}\tau^{2}
&=
c^{2}\mathrm{d}t^{2}
-
|\mathbf{v}|^{2}\mathrm{d}t^{2}\\
&=
c^{2}\mathrm{d}t^{2}
\left(1-
\frac{v^{2}}{c^{2}}
\right).
\end{aligned}
\]

Therefore

\[
\boxed{
\mathrm{d}\tau
=
\mathrm{d}t
\sqrt{1-
\frac{v^{2}}{c^{2}}}
}.
\]

Define the Lorentz factor

\[
\boxed{
\gamma
=
\frac{1}{\sqrt{1-v^{2}/c^{2}}}
}.
\]

Then

\[
\boxed{
\mathrm{d}t
=
\gamma\,\mathrm{d}\tau
}.
\]

Proper time is the time measured by a clock moving along the worldline. It is invariant because it is constructed from the invariant interval.

\subsection{Coordinate time and proper time}

If a clock is at rest in an inertial frame, then

\[
\mathrm{d}\mathbf{x}=0,
\]

so

\[
\mathrm{d}\tau=\mathrm{d}t.
\]

If the clock moves with speed \(v\), then

\[
\mathrm{d}\tau<\mathrm{d}t
\]

for \(v\neq0\). The moving clock accumulates less proper time between the same two events than the coordinate time measured in that frame.

This is time dilation expressed geometrically. The statement follows directly from the interval rather than from a separate dynamical assumption.

\subsection{Worked example: proper time of a moving particle}

A particle moves at constant speed

\[
v=0.80c
\]

for a coordinate time interval

\[
\Delta t=10.0\,\mathrm{s}.
\]

The Lorentz factor is

\[
\gamma
=
\frac{1}{\sqrt{1-0.80^{2}}}
=
\frac{1}{0.60}
=
\frac{5}{3}.
\]

The elapsed proper time is

\[
\Delta\tau
=
\frac{\Delta t}{\gamma}
=
\frac{10.0\,\mathrm{s}}{5/3}
=
6.0\,\mathrm{s}.
\]

Thus the particle's comoving clock records

\[
\boxed{
\Delta\tau=6.0\,\mathrm{s}
}
\]

between the departure and arrival events.

A direct interval calculation gives the same result. Since

\[
\Delta x=v\Delta t=0.80c\Delta t,
\]

we have

\[
\begin{aligned}
\Delta s^{2}
&=
c^{2}(\Delta t)^{2}
-(0.80c\Delta t)^{2}\\
&=
0.36c^{2}(\Delta t)^{2}\\
&=
c^{2}(6.0\,\mathrm{s})^{2}.
\end{aligned}
\]

Hence

\[
\Delta\tau=\frac{\sqrt{\Delta s^{2}}}{c}=6.0\,\mathrm{s}.
\]

\subsection{Worked example: classifying a separation}

Suppose

\[
\Delta t=3.0\,\mu\mathrm{s},
\qquad
|\Delta\mathbf{x}|=1.20\,\mathrm{km}.
\]

The temporal contribution is

\[
c\Delta t
=
(3.00\times10^{8}\,\mathrm{m\,s^{-1}})
(3.0\times10^{-6}\,\mathrm{s})
=
900\,\mathrm{m}.
\]

Therefore

\[
\Delta s^{2}
=
(900\,\mathrm{m})^{2}
-(1200\,\mathrm{m})^{2}.
\]

Calculating,

\[
\Delta s^{2}
=
810000\,\mathrm{m}^{2}
-
1440000\,\mathrm{m}^{2}
=
-630000\,\mathrm{m}^{2}.
\]

Since

\[
\Delta s^{2}<0,
\]

the separation is spacelike.

\subsection{Proper distance for spacelike separation}

For a spacelike separation, one may define the invariant proper distance

\[
\boxed{
\ell_{\mathrm{proper}}
=
\sqrt{-\Delta s^{2}}
}.
\]

The minus sign is needed because \(\Delta s^{2}<0\) with our signature.

There exists an inertial frame in which the two spacelike-separated events are simultaneous:

\[
\Delta t'=0.
\]

In that frame,

\[
\Delta s^{2}
=
-|\Delta\mathbf{x}'|^{2},
\]

so

\[
\ell_{\mathrm{proper}}
=
|\Delta\mathbf{x}'|.
\]

This is the invariant spatial separation measured in the frame where the events are simultaneous.

\subsection{Null intervals}

For light,

\[
\Delta s^{2}=0.
\]

Therefore the proper time along a light ray is

\[
\Delta\tau=0.
\]

This does not mean that a physical rest frame can be attached to light. A massive clock cannot move at speed \(c\), and no inertial frame exists in which a photon is at rest.

The vanishing interval means only that the lightlike path has zero Minkowski length according to the metric.

\subsection{The inverse metric}

The inverse metric \(\eta^{\mu\nu}\) is defined by

\[
\eta^{\mu\alpha}\eta_{\alpha\nu}
=
\delta^{\mu}{}_{\nu},
\]

where

\[
\delta^{\mu}{}_{\nu}
=
\begin{cases}
1,&\mu=\nu,\\
0,&\mu\neq\nu.
\end{cases}
\]

For the Minkowski metric in Cartesian inertial coordinates,

\[
\boxed{
\eta^{\mu\nu}
=
\eta_{\mu\nu}
=
\operatorname{diag}(1,-1,-1,-1)
}.
\]

The equality of the numerical matrices is special to this diagonal convention. The positions of the indices still carry mathematical meaning.

\subsection{Structural checks}

A calculation involving the interval should satisfy:

\begin{enumerate}
    \item every term in \(\Delta s^{2}\) has dimensions of length squared;
    \item the temporal term is positive and spatial terms are negative in the chosen convention;
    \item timelike, null, and spacelike classifications agree with the light-cone picture;
    \item proper time is real only for timelike motion;
    \item the interval remains unchanged after a Lorentz transformation.
\end{enumerate}

\subsection{Common mistakes}

Typical errors include:

\begin{enumerate}
    \item mixing the signatures \((+,-,-,-)\) and \((-,+,+,+)\) in one derivation;
    \item writing the spatial terms with positive signs after choosing the mostly-minus metric;
    \item confusing \(\Delta s\) with \(\Delta s^{2}\);
    \item taking \(\sqrt{\Delta s^{2}}\) for a spacelike interval without accounting for its negative sign;
    \item treating proper time as a coordinate attached to every observer rather than the invariant time along a timelike worldline;
    \item claiming that light has a rest frame because its proper time vanishes;
    \item forgetting the factor \(c\) when comparing time and distance.
\end{enumerate}

\subsection{What has been established}

The Minkowski metric is fixed throughout these notes as

\[
\eta_{\mu\nu}
=
\operatorname{diag}(1,-1,-1,-1).
\]

It defines the invariant interval

\[
\Delta s^{2}
=
c^{2}(\Delta t)^{2}
-
|\Delta\mathbf{x}|^{2},
\]

whose sign classifies separations as timelike, null, or spacelike. For a timelike worldline, the interval defines proper time through

\[
c^{2}\mathrm{d}\tau^{2}=\mathrm{d}s^{2},
\]

leading to

\[
\mathrm{d}t=\gamma\mathrm{d}\tau.
\]

The next section generalizes the coordinate displacement to arbitrary four-vectors and introduces covectors, scalar products, and systematic index raising and lowering.